% ------------------------------------------------------------------------------
\section{Base de datos, copias e intercambio}
\label{sec:bd-base}
% ------------------------------------------------------------------------------

Todos los datos del usuario viven en un único archivo SQLite, \ruta{concretolab.db}, dentro de la carpeta de datos
(capítulo de arquitectura). Este capítulo describe su esquema, cómo se versiona con Flyway, cómo se hacen y restauran
las copias de seguridad, cómo se convierten las bases antiguas protegidas con contraseña, el formato de intercambio de
proyectos, la bitácora, los datos de demostración y la pantalla «Datos y copias».

\subsection{Conexión}

La configuración está en \ruta{src/main/resources/application.properties}:

\begin{table}[H]
	\centering
	\caption{Propiedades de la base de datos y de los registros.}
	\label{tab:bd-propiedades}
	\begin{tabular}{L{7.0cm}L{7.8cm}}
		\hline
		\textbf{Propiedad} & \textbf{Valor y motivo} \\ \hline
		\ruta{spring.datasource.url} & \ruta{jdbc:sqlite:\allowbreak \$\{concretolab.datos:\allowbreak \$\{user.home\}/.concretolab\}/\allowbreak concretolab.db}:
		la carpeta se puede cambiar con \ruta{-Dconcretolab.datos}. \\
		\ruta{spring.datasource.hikari.\allowbreak maximum-pool-size} & \ruta{1}: SQLite admite un solo escritor; con una sola
		conexión no hay bloqueos entre hilos. \\
		\ruta{spring.jpa.database-platform} & \ruta{org.hibernate.community.dialect.\allowbreak{}SQLiteDialect}. \\
		\ruta{spring.jpa.hibernate.ddl-auto} & \ruta{none}: el esquema lo gestiona solo Flyway. \\
		\ruta{spring.jpa.open-in-view} & \ruta{false}. \\
		\ruta{spring.jpa.properties.\allowbreak hibernate.jdbc.time\_zone} & \ruta{UTC}. \\
		\ruta{logging.file.name} & \ruta{<datos>/logs/concretolab.log}; rota a los 5~MB, conserva 7 días y como
		máximo 50~MB. \\ \hline
	\end{tabular}
\end{table}

El controlador JDBC es \ruta{io.github.willena:sqlite-jdbc} 3.46.1.0, que mantiene el paquete \ruta{org.sqlite.JDBC}
e incluye SQLite3MultipleCiphers. La aplicación actual abre la base \textbf{sin cifrar}; el cifrado solo se usa para
reconocer y convertir bases de versiones anteriores (sección~\ref{sec:bd-cifrada}).

\subsection{Tablas y columnas}

El esquema lo crean las migraciones \ruta{V1\_\_esquema\_base.sql} (nueve tablas) y \ruta{V2\_\_bitacora.sql} (tabla
\ruta{bitacora}). Las columnas \ruta{id} son \ruta{integer primary key}, que en SQLite es un alias del \ruta{rowid} y se
autoincrementa; en las entidades son \ruta{Long} con \ruta{GenerationType.IDENTITY}. Los nombres de columna siguen la
estrategia de nombres de Spring (\ruta{proyectoId} $\rightarrow$ \ruta{proyecto\_id}).

\begin{importante}[Sin claves foráneas]
El esquema no declara claves foráneas: las relaciones (\ruta{proyecto\_id}, \ruta{mezcla\_id}, \ruta{material\_id},
\ruta{columna\_id}, \ruta{viga\_id}) son números que mantienen los servicios. Por ejemplo,
\ruta{YefriLermaProyectoService.\allowbreak yefriLermaEliminarProyecto} borra a mano los ensayos, las mezclas, los elementos (con sus
vínculos), la cuadrícula y el expediente antes del proyecto. Quien borre filas por otro camino debe hacer lo mismo.
\end{importante}

\begin{longtable}{L{4.7cm}L{3.1cm}L{6.8cm}}
	\caption{Tablas, columnas y significado.} \label{tab:bd-tablas} \\
	\hline
	\textbf{Columna} & \textbf{Tipo SQL} & \textbf{Significado} \\ \hline
	\endfirsthead
	\hline
	\textbf{Columna} & \textbf{Tipo SQL} & \textbf{Significado} \\ \hline
	\endhead
	\hline
	\endfoot
	\multicolumn{3}{l}{\textbf{\ruta{proyectos}} (\ruta{Proyecto})} \\ \hline
	\ruta{id} & integer, PK & Identificador. \\
	\ruta{nombre} & varchar(120), no nulo & Nombre del proyecto. \\
	\ruta{ubicacion} & varchar(200) & Texto libre. \\ \hline
	\multicolumn{3}{l}{\textbf{\ruta{elementos}} (\ruta{Elemento})} \\ \hline
	\ruta{id} & integer, PK & Identificador. \\
	\ruta{proyecto\_id} & bigint, no nulo & Proyecto dueño. \\
	\ruta{tipo} & varchar(255), no nulo & \ruta{ZAPATA}, \ruta{COLUMNA}, \ruta{VIGA}, \ruta{LOSA} o \ruta{MURO}
	(lo valida el servicio). \\
	\ruta{nombre} & varchar(255), no nulo & Nombre visible (por ejemplo \ruta{C1}). \\
	\ruta{largo}, \ruta{ancho}, \ruta{altura} & float, no nulo & Dimensiones del prisma en X, Y y Z (m). \\
	\ruta{x}, \ruta{y}, \ruta{z} & float, no nulo & Centro del prisma (m). \\ \hline
	\multicolumn{3}{l}{\textbf{\ruta{cuadriculas\_proyecto}} (\ruta{CuadriculaProyecto})} \\ \hline
	\ruta{proyecto\_id} & bigint, PK & Un registro por proyecto. \\
	\ruta{separacionesx}, \ruta{separacionesy} & varchar(2000), no nulo & Separaciones entre ejes, unidas con
	«\ruta{; }». \\
	\ruta{niveles} & varchar(2000), no nulo & Cotas de nivel, unidas con «\ruta{; }». \\
	\ruta{origenx}, \ruta{origeny} & float, no nulo & Origen de la cuadrícula (m). \\ \hline
	\multicolumn{3}{l}{\textbf{\ruta{vinculos\_geometricos}} (\ruta{VinculoGeometrico})} \\ \hline
	\ruta{id} & integer, PK & Identificador. \\
	\ruta{columna\_id}, \ruta{viga\_id} & bigint, no nulo & Columna y viga vinculadas. La entidad declara un
	\ruta{@UniqueConstraint} sobre el par, pero la migración no lo crea. \\ \hline
	\multicolumn{3}{l}{\textbf{\ruta{expedientes\_estructurales}} (\ruta{ExpedienteEstructural})} \\ \hline
	\ruta{proyecto\_id} & bigint, PK & Un expediente por proyecto. \\
	\ruta{ubicacion}, \ruta{uso}, \ruta{fecha\_proyecto}, \ruta{estado\_proyecto}, \ruta{profesional}, \ruta{norma},
	\ruta{justificacion\_norma}, \ruta{estudio\_suelos}, \ruta{observaciones} & varchar(255) & Datos administrativos;
	\ruta{ExpedienteService} exige la fecha como \ruta{AAAA-MM-DD}. \\ \hline
	\multicolumn{3}{l}{\textbf{\ruta{materiales}} (\ruta{Material})} \\ \hline
	\ruta{id} & integer, PK & Identificador. \\
	\ruta{nombre} & varchar(255), no nulo & Nombre del material. \\
	\ruta{tipo} & varchar(255), no nulo & Uno de \ruta{MaterialService.YEFRI\_LERMA\_TIPOS}: \ruta{CEMENTO}, siete
	tipos de cemento (I, IP, ICo, II, V, MS, HS), \ruta{AGREGADO FINO}, \ruta{AGREGADO GRUESO} o \ruta{ADITIVO}. \\
	\ruta{procedencia}, \ruta{referencia\_ensayo}, \ruta{fecha\_ensayo}, \ruta{observaciones} & varchar(255) & Origen
	de los datos. \\
	\ruta{densidad\_relativa}, \ruta{absorcion}, \ruta{humedad}, \ruta{modulo\_finura}, \ruta{tamano\_maximo\_mm},
	\ruta{peso\_unitario\_compactado} & float (admite nulo) & Propiedades medidas (\ruta{Double} en la entidad). \\ \hline
	\multicolumn{3}{l}{\textbf{\ruta{granulometrias}} (\ruta{RegistroGranulometria})} \\ \hline
	\ruta{id} & integer, PK & Identificador. \\
	\ruta{material\_id} & bigint & Agregado ensayado. \\
	\ruta{identificacion}, \ruta{fecha} & varchar(255) & Datos del ensayo. \\
	\ruta{masa\_inicial}, \ruta{fondo} & float, no nulo & Masa seca inicial y masa del fondo (g). \\
	\ruta{mediciones} & varchar(10000) & Una línea por tamiz: \ruta{abertura; masa retenida}. \\ \hline
	\multicolumn{3}{l}{\textbf{\ruta{mezclas}} (\ruta{Mezcla})} \\ \hline
	\ruta{id} & integer, PK & Identificador. \\
	\ruta{proyecto\_id} & bigint & Proyecto de la mezcla. \\
	\ruta{nombre}, \ruta{referencia}, \ruta{observaciones} & varchar(255) & Descripción. \\
	\ruta{material\_fino}, \ruta{material\_grueso} & varchar(255) & Instantánea en texto del agregado elegido (no es
	una referencia a \ruta{materiales}). \\
	\ruta{resistencia} & float, no nulo & f'c especificada (MPa). \\
	\ruta{cemento}, \ruta{agua}, \ruta{fino\_sss}, \ruta{grueso\_sss} & float, no nulo & Cantidades por m$^3$ (kg),
	agregados en condición SSS. \\
	\ruta{absorcion\_fino}, \ruta{humedad\_fino}, \ruta{absorcion\_grueso}, \ruta{humedad\_grueso} & float, no nulo &
	Porcentajes sobre masa seca. \\ \hline
	\multicolumn{3}{l}{\textbf{\ruta{ensayos}} (\ruta{Ensayo})} \\ \hline
	\ruta{id} & integer, PK & Identificador. \\
	\ruta{mezcla\_id} & bigint & Mezcla ensayada. \\
	\ruta{fecha}, \ruta{identificacion}, \ruta{observaciones} & varchar(255) & Datos de la probeta. \\
	\ruta{edad\_dias}, \ruta{resistencia\_mpa}, \ruta{asentamiento\_mm} & float, no nulo & Edad, resistencia a
	compresión y asentamiento. \\ \hline
	\multicolumn{3}{l}{\textbf{\ruta{bitacora}} (\ruta{RegistroBitacora}, migración V2)} \\ \hline
	\ruta{id} & integer, PK & Identificador. \\
	\ruta{fecha} & varchar(30), no nulo & Fecha y hora local ISO-8601 al segundo. \\
	\ruta{accion} & varchar(40), no nulo & Código de la acción (tabla~\ref{tab:bd-acciones}). \\
	\ruta{detalle} & varchar(400) & Descripción legible. \\
\end{longtable}

Flyway agrega su propia tabla, \ruta{flyway\_schema\_history}, con las migraciones aplicadas.

\subsection{Migraciones con Flyway}
\label{sec:bd-flyway}

\begin{table}[H]
	\centering
	\caption{Migraciones existentes.}
	\label{tab:bd-migraciones}
	\begin{tabular}{L{1.4cm}L{5.6cm}L{6.7cm}}
		\hline
		\textbf{Versión} & \textbf{Archivo} & \textbf{Qué hace} \\ \hline
		V1 & \ruta{src/main/resources/db/migration/} \ruta{V1\_\_esquema\_base.sql} & Crea las nueve tablas de las
		entidades con \ruta{CREATE TABLE IF NOT EXISTS} (se generó a partir de las entidades). \\
		V2 & \ruta{src/main/resources/db/migration/} \ruta{V2\_\_bitacora.sql} & Crea la tabla \ruta{bitacora}. \\
		V3 & \ruta{src/main/java/db/migration/} \ruta{V3\_\_Completar\_columnas.java} & Migración Java
		(\ruta{BaseJavaMigration}): lee las sentencias \ruta{CREATE TABLE} de V1, compara con \ruta{PRAGMA table\_info}
		y agrega con \ruta{ALTER TABLE \ldots{} ADD COLUMN} las columnas que falten; a las \ruta{not null} les pone un
		valor por defecto (\ruta{''} o \ruta{0}) para poder agregarlas sobre filas existentes. En una base nueva no hace
		nada. \\ \hline
	\end{tabular}
\end{table}

Flyway busca en su ubicación por defecto, \ruta{classpath:db/migration}, que abarca tanto los \ruta{.sql} de
\ruta{resources/db/migration} como las clases Java del paquete \ruta{db.migration}.

\begin{nota}[V3 sí existe, pero en otra carpeta]
La migración V3 existe y se compila (\ruta{target/classes/db/migration/V3\_\_Completar\_columnas.class}), pero es una
clase Java en \ruta{src/main/java/db/migration}, no un \ruta{.sql}. Por eso no aparece al listar
\ruta{src/main/resources/db/migration}. El README la menciona («V3 columnas faltantes») junto a esa carpeta de recursos;
el comentario de \ruta{application.properties} solo nombra V1 y V2. \ruta{EsquemaMigracionesTest} comprueba en
\ruta{flyway\_schema\_history} que se aplicaron las versiones 1 y 2, pero no la 3. V3 se ejercita de forma indirecta en
\ruta{CompatibilidadMezclaTest}: crea una tabla \ruta{mezclas} sin \ruta{material\_fino} ni \ruta{material\_grueso} y
comprueba que la entidad \ruta{Mezcla} se puede leer después de migrar.
\end{nota}

Las propiedades de Flyway son:

\begin{itemize}
	\item \ruta{spring.flyway.enabled=true}.
	\item \ruta{spring.flyway.baseline-on-migrate=true} y \ruta{spring.flyway.baseline-version=0}. Las primeras versiones
	      de ConcretoLab creaban el esquema con \ruta{ddl-auto=update} y no tienen la tabla
	      \ruta{flyway\_schema\_history}. Flyway se niega a migrar una base no vacía sin historial; con estas opciones la
	      toma como punto de partida (versión 0) y aplica V1, V2 y V3 encima. V1 no falla porque usa
	      \ruta{IF NOT EXISTS}, y V3 completa las columnas que esas bases no tuvieran.
	\item \ruta{spring.jpa.hibernate.ddl-auto=none}. Hibernate ya no modifica el esquema: \textbf{cada cambio en una
	      entidad necesita una migración nueva}. No se usa \ruta{validate} porque Hibernate espera \ruta{bigint} en los
	      \ruta{id}, mientras que SQLite necesita \ruta{integer} para que la clave primaria se autoincremente; la validación
	      fallaría siempre. En su lugar, \ruta{EsquemaMigracionesTest} consulta cada entidad sobre una base migrada (nueva
	      y «antigua») y falla si falta una tabla o una columna.
\end{itemize}

\begin{advertencia}[Antes de migrar se guarda una copia]
\ruta{JavaFxApplication.init()} guarda la copia «inicio» \textbf{antes} de levantar Spring, es decir, antes de que
Flyway toque la base. Si una migración nueva fallara en la computadora de un cliente, esa copia conserva los datos tal
como estaban.
\end{advertencia}

\subsection{Copias de seguridad}
\label{sec:bd-copias}

Hay dos clases complementarias:

\begin{itemize}
	\item \ruta{RespaldoDatos}: utilidad \textbf{sin Spring} (solo JDK y \ruta{System.Logger}), porque se usa en
	      \ruta{init()} y \ruta{stop()}, cuando el contexto todavía no existe o ya se cerró. Su regla es que ninguna
	      operación de arranque lanza excepciones: una copia que falla se registra y no impide abrir la aplicación.
	\item \ruta{RespaldoService}: bean de Spring para la base \textbf{abierta} (copias manuales, listado, eliminación y
	      programación de la restauración).
\end{itemize}

\subsubsection{Nombres y rotación}

Cada copia se llama \ruta{concretolab-AAAAMMDD-HHMMSS-<motivo>.db} y se guarda en \ruta{<datos>/backups}. El patrón
\ruta{ERICK\_DELGADO\_PATRON\_NOMBRE} reconoce la marca de tiempo y el motivo; los archivos que no lo cumplen se
ignoran al listar. El motivo se normaliza (minúsculas, sin tildes, solo letras y guiones; vacío = «manual»). Si ya
existe una copia con la misma marca, se avanza un segundo hasta encontrar un nombre libre.

\begin{table}[H]
	\centering
	\caption{Motivos de copia, cuándo se crean y cuántas se conservan.}
	\label{tab:bd-motivos}
	\begin{tabular}{L{3.7cm}L{2.8cm}L{4.4cm}L{2.2cm}}
		\hline
		\textbf{Motivo} & \textbf{Texto en pantalla} & \textbf{Cuándo} & \textbf{Se conservan} \\ \hline
		\ruta{inicio} & Al abrir la aplicación & \ruta{init()}, antes de Spring y Flyway & 10 \\
		\ruta{cierre} & Al cerrarla & \ruta{stop()}, con el contexto ya cerrado & 10 \\
		\ruta{manual} & Manual & Botón «Crear copia ahora» & Todas \\
		\ruta{previa-restauracion} & Antes de restaurar & Al aplicar una restauración & 20 \\
		\ruta{antes-de-quitar-clave} & Base anterior con contraseña & Al convertir una base antigua & Todas \\ \hline
	\end{tabular}
\end{table}

La rotación (\ruta{erickDelgadoRotar}) ordena las copias de un motivo por nombre, de la más reciente a la más antigua
(la marca de tiempo está en el nombre), y borra las que exceden el límite. Solo se rotan \ruta{inicio}, \ruta{cierre}
y \ruta{previa-restauracion}.

\subsubsection{Cómo se copia}

\begin{itemize}
	\item \textbf{Copias de arranque y cierre} (\ruta{RespaldoDatos.\allowbreak{}erickDelgadoRespaldarArchivo}): copian el archivo
	      \ruta{.db} cuando la base no está abierta. Primero se copia a un \ruta{.tmp} en la misma carpeta y luego se
	      renombra, de modo que nunca queda una copia a medias con un nombre válido. Si la base no existe o está vacía, no
	      hacen nada.
	\item \textbf{Copias manuales} (\ruta{RespaldoService.\allowbreak erickDelgadoCrearCopia}): la base está abierta, así que se
	      usa \ruta{VACUUM INTO '<destino>'}, que escribe un archivo coherente aunque haya conexiones activas. El método
	      no es transaccional (\ruta{VACUUM} no puede ejecutarse dentro de una transacción) y libera la conexión antes de
	      escribir en la bitácora, porque el pool tiene una sola. Si falla, borra el destino incompleto.
\end{itemize}

\subsubsection{Restauración diferida}

La base no se reemplaza en caliente. El proceso tiene dos fases:

\begin{enumerate}
	\item \textbf{Programar} (\ruta{erickDelgadoProgramarRestauracion}): valida la copia y escribe su ruta absoluta en
	      el marcador \ruta{<datos>/restaurar.pendiente}. La pantalla lo registra en la bitácora
	      (\ruta{RESTAURACION\_PROGRAMADA}) y cierra la aplicación con \ruta{Platform.exit()}.
	\item \textbf{Aplicar} (\ruta{erickDelgadoAplicarRestauracionPendiente}, primera instrucción de \ruta{init()} después
	      de crear la carpeta): lee el marcador, comprueba que la copia existe y es válida, guarda la base actual como
	      \ruta{previa-restauracion} (si esa copia falla, cancela), copia la elegida a \ruta{concretolab.db.restaurando},
	      borra los archivos \ruta{-wal}, \ruta{-shm} y \ruta{-journal} y la mueve sobre \ruta{concretolab.db}. El marcador
	      se borra siempre, haya salido bien o mal, para no entrar en un bucle.
\end{enumerate}

La validación (\ruta{erickDelgadoEsSqliteValido}) acepta un archivo que empieza con la cabecera
\ruta{SQLite format 3\textbackslash{}0}. También acepta uno sin esa cabecera si mide al menos 4~KB y es múltiplo de 512
bytes: es el aspecto de una base cifrada de una versión anterior, que luego se convierte en el arranque.

\begin{consejo}[Eliminar copias]
\ruta{erickDelgadoEliminarCopia} solo borra archivos situados directamente en la carpeta \ruta{backups}; cualquier otra
ruta lanza \ruta{IllegalArgumentException}. La eliminación no se anota en la bitácora.
\end{consejo}

\subsection{Bases antiguas protegidas con contraseña}
\label{sec:bd-cifrada}

Las versiones anteriores de ConcretoLab cifraban la base (AES-256, modo SQLCipher~4 de SQLite3MultipleCiphers). La
actual abre sin contraseña y convierte esas bases \textbf{una sola vez}. Intervienen \ruta{BaseCifrada},
\ruta{DialogoClaveBase} y \ruta{JavaFxApplication}.

\begin{enumerate}
	\item \textbf{Detección.} \ruta{BaseCifrada.\allowbreak{}erickDelgadoInspeccionar} lee los primeros 16 bytes:
	      \ruta{NO\_EXISTE} (no hay archivo o está vacío), \ruta{SIN\_CIFRAR} (cabecera SQLite) o \ruta{CIFRADA}
	      (cualquier otra cosa).
	\item \textbf{Contraseña.} Si está cifrada, \ruta{JavaFxApplication} pide la contraseña con
	      \ruta{DialogoClaveBase.\allowbreak{}erickDelgadoPedirActual}, hasta 3 intentos. En el primer intento usa la variable de
	      entorno \ruta{CONCRETOLAB\_DB\_CLAVE} si existe (para automatizar). La contraseña se comprueba leyendo
	      \ruta{sqlite\_master}; si es incorrecta, no se modifica nada. Si la persona pulsa «Salir» o se agotan los
	      intentos, la aplicación se cierra sin tocar la base.
	\item \textbf{Conversión} (\ruta{BaseCifrada.\allowbreak{}erickDelgadoQuitarClave}):
	      \begin{enumerate}
		      \item cuenta las filas de cada tabla con la contraseña;
		      \item copia el archivo cifrado a \ruta{backups} con el motivo \ruta{antes-de-quitar-clave} y comprueba que la
		            copia tiene el mismo tamaño y sigue cifrada (si no, aborta sin tocar la base);
		      \item ejecuta \ruta{PRAGMA rekey=''} para quitar la contraseña;
		      \item verifica que la cabecera ahora es SQLite, que \ruta{PRAGMA integrity\_check} devuelve \ruta{ok} y que
		            cada tabla tiene las mismas filas que antes;
		      \item si cualquier paso falla, \textbf{vuelve a copiar la base cifrada en su lugar} y la deja como estaba.
	      \end{enumerate}
	\item \textbf{Arranque.} Luego se guarda la copia «inicio» y Spring abre la base con
	      \ruta{BaseCifrada.\allowbreak{}erickDelgadoPropiedadesSpring(null)} (\ruta{cipher=sqlcipher}, \ruta{legacy=4}, sin clave,
	      es decir, sin cifrar). El arreglo \ruta{char[]} con la contraseña se llena de ceros al terminar.
\end{enumerate}

\ruta{BaseCifrada} conserva además \ruta{erickDelgadoCifrarExistente} y \ruta{erickDelgadoCambiarClave}, de cuando la
aplicación cifraba la base. Ninguna pantalla las usa; solo las ejercita \ruta{CifradoBaseTest}.

\subsection{Exportar e importar proyectos}
\label{sec:bd-intercambio}

\ruta{IntercambioProyectoService} lleva un proyecto de una computadora a otra con un archivo \ruta{.concretolab}.

\subsubsection{Formato}

Un \ruta{.concretolab} es un ZIP con exactamente dos entradas:

\begin{itemize}
	\item \ruta{manifiesto.json}: \ruta{formato} (\ruta{concretolab-proyecto}), \ruta{version} (1),
	      \ruta{aplicacion}, \ruta{fecha} (instante ISO), \ruta{sha256} de \ruta{datos.json} y \ruta{conteos}
	      (elementos, vínculos, mezclas y ensayos).
	\item \ruta{datos.json}: el proyecto (nombre y ubicación), sus elementos, la cuadrícula, los vínculos (referidos por
	      la \textbf{posición} del elemento en la lista, no por su \ruta{id}), las mezclas con sus ensayos y el expediente.
\end{itemize}

No viajan la biblioteca de materiales ni los tamizados, porque no pertenecen a un proyecto; las mezclas conservan como
texto el nombre de sus agregados. El archivo se escribe de forma atómica: primero un temporal en la misma carpeta y
luego \ruta{Files.move} con \ruta{ATOMIC\_MOVE} (o un movimiento normal si el sistema de archivos no lo admite).

\subsubsection{Cifrado opcional del archivo}

Si al exportar se marca «Proteger este archivo con contraseña», el ZIP completo se cifra con \ruta{CifradoArchivo}.
\textbf{Solo afecta al archivo exportado}, no a la base local.

\begin{table}[H]
	\centering
	\caption{Formato de un archivo cifrado por \ruta{CifradoArchivo}.}
	\label{tab:bd-cifrado}
	\begin{tabular}{L{4.4cm}L{9.4cm}}
		\hline
		\textbf{Parte} & \textbf{Detalle} \\ \hline
		Cabecera & \ruta{CLCIF1} (6 bytes ASCII) + sal (16 bytes) + IV (12 bytes). \\
		Clave & PBKDF2WithHmacSHA256, 210\,000 iteraciones, 256 bits, con la sal. \\
		Cuerpo & AES/GCM/NoPadding con etiqueta de 16 bytes; la cabecera se autentica como datos adicionales (AAD), así
		que alterarla también se detecta. \\
		Errores & Contraseña incorrecta, archivo truncado o modificado dan el mismo mensaje genérico: «Contraseña
		incorrecta o archivo dañado.» \\ \hline
	\end{tabular}
\end{table}

\ruta{CifradoArchivo} solo exige que la contraseña no esté vacía; la pantalla «Datos y copias» pide al menos 8
caracteres y una confirmación al exportar.

\subsubsection{Importación}

\ruta{erickDelgadoImportar} es \ruta{@Transactional} y valida todo antes de guardar:

\begin{enumerate}
	\item el archivo no supera 50~MB; si empieza con \ruta{CLCIF1}, se descifra con la contraseña;
	\item el ZIP solo puede contener \ruta{manifiesto.json} (hasta 1~MB) y \ruta{datos.json} (hasta 100~MB), una vez cada
	      uno; se rechazan rutas con \ruta{..}, \ruta{/}, \ruta{\textbackslash}, \ruta{:} y directorios;
	\item el manifiesto debe tener el formato correcto y una versión no mayor que la soportada (si es mayor, se pide
	      actualizar ConcretoLab); el SHA-256 de \ruta{datos.json} debe coincidir (comparación de tiempo constante);
	\item los datos se validan: hasta 20\,000 elementos, 2\,000 mezclas y 20\,000 ensayos; tipos válidos; dimensiones
	      positivas y finitas; vínculos que unan una columna con una viga existentes (los repetidos se omiten con aviso);
	      textos acotados; humedades y absorciones entre 0 y 100~\%; fechas de ensayo ISO;
	\item se crea \textbf{siempre un proyecto nuevo}, nunca se reemplaza uno existente. Si el nombre ya existe se agrega
	      « (importado)», « (importado 2)», etc., sin pasar de 120 caracteres, y se devuelve un aviso.
\end{enumerate}

Si algo falla durante el guardado, la transacción revierte todo y no queda un proyecto a medias.

\subsection{Bitácora de cambios}

\ruta{BitacoraService.\allowbreak{}erickDelgadoRegistrar(accion, detalle)} guarda un \ruta{RegistroBitacora}. Es
\ruta{@Transactional} y se une a la transacción de la operación que registra (SQLite trabaja con una sola conexión y no
puede abrir otra), así que el evento se confirma o se revierte junto con ella. Recorta la acción a 40 caracteres y el
detalle a 400 (terminando en «\ldots»). Si no puede guardar, lo anota en el registro técnico y no interrumpe la
operación. \ruta{erickDelgadoRecientes()} devuelve los últimos 200 eventos.

\begin{table}[H]
	\centering
	\caption{Acciones que se registran hoy.}
	\label{tab:bd-acciones}
	\begin{tabular}{L{5.0cm}L{8.8cm}}
		\hline
		\textbf{Acción} & \textbf{Dónde se registra} \\ \hline
		\ruta{PROYECTO\_CREADO} & \ruta{YefriLermaProyectoService.\allowbreak yefriLermaCrearProyecto} \\
		\ruta{PROYECTO\_ELIMINADO} & \ruta{YefriLermaProyectoService.\allowbreak yefriLermaEliminarProyecto} \\
		\ruta{PROYECTO\_EXPORTADO} & \ruta{IntercambioProyectoService.\allowbreak erickDelgadoExportar} (indica si lleva contraseña) \\
		\ruta{PROYECTO\_IMPORTADO} & \ruta{IntercambioProyectoService.\allowbreak erickDelgadoImportar} (con los conteos) \\
		\ruta{COPIA\_CREADA} & \ruta{RespaldoService.\allowbreak erickDelgadoCrearCopia} (solo copias manuales) \\
		\ruta{RESTAURACION\_PROGRAMADA} & \ruta{RespaldoService.\allowbreak erickDelgadoProgramarRestauracion} \\ \hline
	\end{tabular}
\end{table}

No se registran las ediciones de elementos, mezclas o materiales, la eliminación de copias, la restauración ya
aplicada ni la conversión de una base con contraseña (esta última queda en el registro técnico y en la copia
\ruta{antes-de-quitar-clave}).

\subsection{Datos de demostración (seeders)}

Los seeders del paquete \ruta{seeder} crean proyectos y datos ficticios para aprender y probar.
\ruta{DatabaseSeeder} es un \ruta{ApplicationRunner} anotado con
\ruta{@ConditionalOnProperty} sobre la propiedad \ruta{app.seeder.enabled} con valor \ruta{true}: si la propiedad no vale
\ruta{true}, el bean ni siquiera existe. En \ruta{application.properties} está en \ruta{false}.

\begin{advertencia}[Apagados por defecto]
Al activar los seeders sobre la carpeta de datos real, los proyectos de demostración se crean en la base del usuario,
y si se borran vuelven a aparecer en cada arranque mientras la opción siga activa. Actívelos solo con una carpeta de
prueba (\ruta{-Dconcretolab.datos}).
\end{advertencia}

Formas de activarlos: \ruta{app.seeder.enabled=true} en \ruta{application.properties},
\ruta{java -Dapp.seeder.enabled=true -jar target/concretolab-desktop-0.1.0.jar}, o
\ruta{CREAR\_DISTRIBUCION\_WINDOWS.ps1 -ConDatosDemostracion}, que agrega esa opción al lanzador generado.

\ruta{DatabaseSeeder.\allowbreak{}fernandoHuancaSembrar()} los ejecuta en orden de dependencias y devuelve cuántos registros creó.
Todos son idempotentes: buscan por nombre y no duplican.

\begin{table}[H]
	\centering
	\caption{Datos que crea cada seeder, en orden de ejecución.}
	\label{tab:bd-seeders}
	\begin{tabular}{L{4.4cm}L{9.4cm}}
		\hline
		\textbf{Seeder} & \textbf{Crea} \\ \hline
		\ruta{ProyectoSeeder} & «DEMO - Edificación de práctica» con la plantilla rectangular (4 zapatas, 4 columnas,
		4 vigas). \\
		\ruta{MaterialSeeder} & Una arena, una piedra y un cemento (omite los que ya existen por nombre). \\
		\ruta{GranulometriaSeeder} & Un tamizado de la arena DEMO. \\
		\ruta{MezclaSeeder} & Una dosificación de 21~MPa con un ensayo a 28 días en el proyecto DEMO. \\
		\ruta{EdificiosSismicosSeeder} & «DEMO - Pórtico regular de 3 pisos», «DEMO - Edificio asimétrico de 4 pisos» y
		«DEMO - Torre flexible de 6 pisos» (cumple, torsión y deriva excesiva). \\
		\ruta{EdificioPracticaSeeder} & «Práctica - Edificio de dos pisos incompleto», con cuadrícula y elementos
		faltantes para completar. \\ \hline
	\end{tabular}
\end{table}

\subsection{La pantalla «Datos y copias»}

La vista \ruta{resources/view/datos-view.fxml} y su controlador \ruta{DatosController} están pensados para el
cliente. Tiene tres bloques y un panel lateral:

\begin{longtable}{L{3.8cm}L{10.0cm}}
	\caption{Bloques de la pantalla «Datos y copias».} \label{tab:bd-pantalla} \\
	\hline
	\textbf{Bloque} & \textbf{Controles y comportamiento} \\ \hline
	\endfirsthead
	\hline
	\textbf{Bloque} & \textbf{Controles y comportamiento} \\ \hline
	\endhead
	\hline
	\endfoot
	Copias de seguridad & «Crear copia ahora» (\ruta{onCrearCopia}, motivo \ruta{manual}); tabla Fecha, Tipo y
	Tamaño (\ruta{tablaCopias}); «Restaurar la seleccionada\ldots» pide confirmación con
	\ruta{Notificador.\allowbreak{}fernandoHuancaConfirmar}, programa la restauración y cierra la aplicación; «Eliminar copia» pide
	confirmación. Un aviso (\ruta{lblRestauracionPendiente}) muestra si hay una restauración programada. \\
	Llevar un proyecto a otra computadora & Selector de proyecto (\ruta{cmbProyecto}, preseleccionado con el de
	\ruta{SesionProyecto}), «Exportar\ldots» e «Importar\ldots». Al importar, el proyecto nuevo pasa a ser el activo. \\
	Opciones avanzadas & Panel plegado (\ruta{panelAvanzado}): rutas de la base, de las copias y de los registros,
	tamaño de la base, «Abrir carpeta de datos», «Abrir carpeta de registros» y el historial de cambios (los 200
	eventos más recientes de la bitácora, con «Actualizar»). \\
	Panel de contraseña & \ruta{panelContrasena}, abierto con \ruta{Notificador.\allowbreak{}fernandoHuancaAbrirPanel}. Al exportar:
	casilla «Proteger este archivo con contraseña (opcional)», contraseña de al menos 8 caracteres y confirmación; luego
	se elige el destino. Al importar un archivo cifrado: solo la contraseña. Los campos se vacían al terminar y los
	\ruta{char[]} se borran. \\
\end{longtable}

Exportar, importar y crear copias se ejecutan fuera del hilo de JavaFX, en un \ruta{Task} sobre un hilo \emph{daemon}
llamado \ruta{datos} (\ruta{erickDelgadoEnSegundoPlano}). Los \ruta{IllegalArgumentException} se muestran tal cual,
porque sus mensajes están escritos para el usuario; los demás errores se anotan en el registro técnico y se muestran con
un prefijo genérico. La acción de cerrar la aplicación está en un campo (\ruta{erickDelgadoCerrarAplicacion}, por
defecto \ruta{Platform::exit}).
